% ------------------------------------------------------------------
%  Holder Gibbs finite-cylinder holes: block recoding and pressure
% ------------------------------------------------------------------

\section{H\"older cylinder holes after block recoding}
\label{app:holder-finite-cylinder}

This appendix is independent of the finite-state resolvent theorem.  It
has two limited purposes.  First, for a one-step cylinder hole it records
the actual H\"older transfer-operator statement on the Banach space
\(C^\alpha\): the open operator is defined explicitly, its leading
eigenvalue and adjoint eigenfunctional are obtained from the survivor
Ruelle--Perron--Frobenius theorem, and the complementary part is shown to
have smaller spectral radius.  Second, after a finite block recoding, it
proves the finite-prefix Bayes-error comparison needed for general Gibbs
weights.  No statement in this appendix is a hypothesis for the
finite-state cyclotomic resolvent theorem.

The Bayes-error comparison has four finite steps.  First, Gibbs
distortion gives uniform lower bounds for bounded-length continuations.
Second, a finite two-terminal reachability condition gives uniform
conditional ambiguity on every surviving prefix.  Third, a safe-prefix
decomposition identifies the Bayes error with a survivor-cylinder sum.
Fourth, the standard pressure formula for the closed survivor subshift
gives the exponential rate of that sum.

We use the following standard form of the Gibbs property.  If \(Z\) is a
topologically mixing one-sided subshift of finite type, \(\psi\) is
H\"older, and \(\nu_\psi\) is its equilibrium state, then there are
constants \(C_G\geq1\) and \(V_\psi<\infty\) such that for every
admissible word \(w=w_0\cdots w_{n-1}\) and every \(x\in[w]\),
\[
  C_G^{-1}e^{S_n\psi(x)-nP_Z(\psi)}
  \leq \nu_\psi([w])
  \leq
  C_G e^{S_n\psi(x)-nP_Z(\psi)},
\]
and \(\abs{S_n\psi(x)-S_n\psi(y)}\leq V_\psi\) whenever
\(x,y\in[w]\).

\begin{lemma}[Uniform Gibbs continuation]
\label{lem:holder-gibbs-continuation}
Let \(Z,\psi,\nu_\psi\) be as above.  For every integer \(L\geq1\) there
is a constant \(c_L>0\) such that, whenever \(w\) and \(v\) are
admissible words, \(\abs{v}\leq L\), and the concatenation \(wv\) is
admissible,
\[
  \frac{\nu_\psi([wv])}{\nu_\psi([w])}\geq c_L .
\]
\end{lemma}

\begin{proof}
Write \(n=\abs w\) and \(k=\abs v\leq L\), and choose
\(x\in[wv]\).  The Gibbs estimate applied to \(wv\), and then to \(w\),
gives
\[
  \frac{\nu_\psi([wv])}{\nu_\psi([w])}
  \geq
  C_G^{-2}
  \exp\bigl(S_{n+k}\psi(x)-(n+k)P_Z(\psi)
            -\sup_{y\in[w]}(S_n\psi(y)-nP_Z(\psi))\bigr).
\]
Since \(x\in[w]\), the Bowen distortion bound gives
\[
  S_n\psi(x)-\sup_{y\in[w]}S_n\psi(y)\geq -V_\psi.
\]
The remaining \(k\) iterates are bounded below by
\(-L\norm{\psi}_\infty\), while \(-kP_Z(\psi)\geq -L\abs{P_Z(\psi)}\).
Thus
\[
  \frac{\nu_\psi([wv])}{\nu_\psi([w])}
  \geq
  C_G^{-2}
  \exp\bigl(-V_\psi-L\norm{\psi}_\infty
              -L\abs{P_Z(\psi)}\bigr),
\]
which is positive and depends only on \(L\), \(Z\), and \(\psi\).
\end{proof}

\begin{lemma}[Uniform finite-prefix ambiguity]
\label{lem:holder-extension-ambiguity}
Let \(Z\) be a topologically mixing one-sided subshift of finite type,
let \(\psi\) be a H\"older potential, and let \(\nu_\psi\) be its
equilibrium state.  Let \(H\) be a non-empty proper one-step Markov hole
with safe alphabet \(\mathcal S\).  Fix \(q\geq2\) and
\(\emptyset\neq R\subsetneq\ZZ/q\ZZ\).  Assume the finite two-terminal
residue reachability condition on the graph
\[
  \mathcal S\times\ZZ/q\ZZ,\qquad \star_R,\quad\star_{R^c},
\]
where safe edges increase the residue by one along admissible safe
transitions, and a terminal edge records the residue of the first
admissible transition into \(H\): both terminal vertices are reachable
from every safe-residue vertex.

Let
\[
  \tau_H(z):=\inf\{n\geq0:z_n\notin\mathcal S\},
  \qquad
  \mathsf P_R:=\{z:\tau_H(z)<\infty,\ \tau_H(z)\in R\bmod q\}.
\]
Then there is \(c_0>0\) such that, for every admissible safe word
\(w=w_0\cdots w_{n-1}\) of length \(n\geq1\),
\[
  c_0
  \leq
  \nu_\psi(\mathsf P_R\mid[w])
  \leq
  1-c_0 .
\]
\end{lemma}

\begin{proof}
For each pair \((i,r)\in\mathcal S\times\ZZ/q\ZZ\), the reachability
condition gives two admissible first-hit continuations from \(i\): one
whose first hit has residue in \(R-r\), and one whose first hit has
residue in \((\ZZ/q\ZZ\setminus R)-r\).  Concretely, these are finite
words \(v^{R}_{i,r}\) and \(v^{R^c}_{i,r}\), appended after the symbol
\(i\), whose intermediate symbols are safe, whose last symbol lies in
the hole, and whose first-hit residues force \(\mathsf P_R\) and
\(\mathsf P_R^c\), respectively.  There are only finitely many
\((i,r)\), so let \(L\) be the maximum of the two continuation lengths.

Let \(w=w_0\cdots w_{n-1}\) be safe, put \(i=w_{n-1}\), and put
\(r=n-1\bmod q\).  The cylinders \([wv^{R}_{i,r}]\) and
\([wv^{R^c}_{i,r}]\) are contained in \([w]\), and they force
\(\mathsf P_R\) and \(\mathsf P_R^c\), respectively.  Lemma
\ref{lem:holder-gibbs-continuation} gives a constant \(c_L>0\) such that
\[
  \nu_\psi([wv^{R}_{i,r}]\mid[w])\geq c_L,
  \qquad
  \nu_\psi([wv^{R^c}_{i,r}]\mid[w])\geq c_L .
\]
Taking \(c_0=\min\{c_L,\tfrac12\}\) gives
\[
  \nu_\psi(\mathsf P_R\mid[w])\geq c_0,
  \qquad
  \nu_\psi(\mathsf P_R^c\mid[w])\geq c_0,
\]
which is the claimed two-sided ambiguity bound.
\end{proof}

\begin{lemma}[Safe-prefix Bayes decomposition]
\label{lem:holder-safe-prefix-bayes-decomposition}
Let \(Z\) be a topologically mixing one-sided subshift of finite type,
let \(\nu\) be a fully supported probability measure on its cylinders,
and let \(H\) be a one-step Markov hole with safe alphabet
\(\mathcal S\).  Fix \(q\geq2\), a non-empty proper
\(R\subsetneq\ZZ/q\ZZ\), and
\[
  \tau_H(z):=\inf\{n\geq0:z_n\notin\mathcal S\},
  \qquad
  \mathsf P_R:=\{z:\tau_H(z)<\infty,\ \tau_H(z)\in R\bmod q\}.
\]
Then for every \(n\geq1\),
\[
  \varepsilon_n(\mathsf P_R;\nu)
  =
  \sum_{\substack{w\in\mathcal L_n(Z)\\
                  w_0,\ldots,w_{n-1}\in\mathcal S}}
    \nu([w])\,
    \min\{
      \nu(\mathsf P_R\mid[w]),
      1-\nu(\mathsf P_R\mid[w])
    \}.
\]
\end{lemma}

\begin{proof}
Apply the cylinder Bayes formula
Proposition~\ref{prop:scan-error-cylinder} to the length-\(n\) cylinder
partition.  If \(w=w_0\cdots w_{n-1}\) contains a symbol outside
\(\mathcal S\), let \(t<n\) be its first such position.  Then
\(\tau_H=t\) on the whole cylinder \([w]\), so the truth value of
\(\mathsf P_R\) is already known from the prefix and the corresponding
Bayes factor is zero.  The remaining cylinders are exactly the safe
prefixes in the displayed sum.
\end{proof}

\begin{lemma}[Safe-word pressure]
\label{lem:holder-safe-word-pressure}
Let \(Z,\psi,H,\mathcal S\) be as in
Lemma~\ref{lem:holder-extension-ambiguity}, and assume that the survivor
subshift
\[
  Z_H:=\{z\in Z:z_n\in\mathcal S\text{ for all }n\geq0\}
\]
is topologically mixing.  Then
\[
  \lim_{n\to\infty}
  -\frac1n\log \nu_\psi(\tau_H\geq n)
  =
  P_Z(\psi)-P_{Z_H}(\psi).
\]
\end{lemma}

\begin{proof}
Let \(\mathcal L_n^{\mathrm{safe}}\) be the admissible length-\(n\)
words all of whose symbols lie in \(\mathcal S\), and set
\[
  Z_n^{\mathrm{safe}}(\psi)
  :=
  \sum_{w\in\mathcal L_n^{\mathrm{safe}}}
    \exp\bigl(\sup_{[w]}S_n\psi\bigr).
\]
The Gibbs estimate gives constants \(C_1,C_2>0\) such that
\[
  C_1e^{-nP_Z(\psi)}Z_n^{\mathrm{safe}}(\psi)
  \leq
  \nu_\psi(\tau_H\geq n)
  \leq
  C_2e^{-nP_Z(\psi)}Z_n^{\mathrm{safe}}(\psi).
\]
It remains to compute the exponential growth of
\(Z_n^{\mathrm{safe}}(\psi)\).

Let \(\mathcal S_\infty\subseteq\mathcal S\) be the symbols that occur
in \(Z_H\), and put \(\mathcal T=\mathcal S\setminus\mathcal S_\infty\).
There is no safe edge from \(\mathcal T\) into \(\mathcal S_\infty\);
otherwise an infinite survivor path beginning in \(\mathcal S_\infty\)
could be appended after that edge, placing the initial transient symbol
in \(Z_H\).  The directed graph on \(\mathcal T\) is acyclic, since a
cycle in \(\mathcal T\) could be repeated indefinitely to produce a
survivor orbit through a transient symbol.  Hence there is a number
\(K\) such that every safe word is a survivor word over
\(\mathcal S_\infty\), possibly followed by a transient suffix of length
at most \(K\).  Words of total length at most \(K\) are irrelevant for
the limit.

Let
\[
  Z_n^{H}(\psi)
  :=
  \sum_{\substack{w\text{ admissible in }Z,\ \abs w=n\\
                  w_0,\ldots,w_{n-1}\in\mathcal S_\infty}}
    \exp\bigl(\sup_{[w]\cap Z_H}S_n\psi\bigr)
\]
be the usual survivor partition sum.  Since survivor words are safe,
\(Z_n^{\mathrm{safe}}(\psi)\geq Z_n^H(\psi)\).  Conversely, if
\(w=uv\) with \(u\) a survivor word and \(v\) a transient suffix of
length \(k\leq K\), then
\[
  \exp\bigl(\sup_{[w]}S_n\psi\bigr)
  \leq
  e^{V_\psi+K\norm{\psi}_\infty}
  \exp\bigl(\sup_{[u]\cap Z_H}S_{n-k}\psi\bigr).
\]
Here the factor \(e^{V_\psi}\) compares the ambient supremum over the
survivor prefix cylinder \([u]\) with the supremum over
\([u]\cap Z_H\); the latter is non-empty because \(u\) is a word in the
mixing survivor subshift.
There are only finitely many possible transient suffixes of length at
most \(K\).  Therefore, for a constant \(C_3<\infty\),
\[
  Z_n^{\mathrm{safe}}(\psi)
  \leq
  C_3\sum_{k=0}^{K} Z_{n-k}^{H}(\psi)
\]
for all large \(n\).  Because \(Z_H\) is a mixing subshift of finite type
and \(\psi|_{Z_H}\) is H\"older,
\[
  \lim_{n\to\infty}\frac1n\log Z_n^H(\psi)
  =
  P_{Z_H}(\psi)
\]
by the standard pressure formula for H\"older potentials on mixing SFTs
\citep{Bowen2008,Ruelle1978}.  The preceding lower and upper bounds show
that \(Z_n^{\mathrm{safe}}(\psi)\) has the same exponential growth rate.
Combining this with the Gibbs comparison for
\(\nu_\psi(\tau_H\geq n)\) proves the displayed escape exponent.
\end{proof}

\begin{theorem}[Open transfer operator on the H\"older space]
\label{thm:holder-open-transfer}
Let \(Z,\psi,H,\mathcal S\) be as in
Lemma~\ref{lem:holder-extension-ambiguity}, and assume that the survivor
subshift
\[
  Z_H=\{z\in Z:z_n\in\mathcal S\text{ for all }n\geq0\}
\]
is topologically mixing.  Choose a standard symbolic metric
\(d_\theta\), \(0<\theta<1\), for which \(\psi\) is Lipschitz, and set
\[
  \mathcal B_\theta(Z):=
  \{f:Z\to\mathbb C:\norm{f}_\theta<\infty\},
  \qquad
  \norm{f}_\theta:=\norm{f}_\infty+
  \sup_{x\neq y}\frac{\abs{f(x)-f(y)}}{d_\theta(x,y)} .
\]
Define the open Ruelle operator
\[
  (\mathcal L_{\psi,H}f)(x)
  :=
  \sum_{\substack{\sigma y=x\\ y_0\in\mathcal S}}
    e^{\psi(y)}f(y),
  \qquad f\in\mathcal B_\theta(Z).
\]
Then \(\mathcal L_{\psi,H}\) is a bounded operator on
\(\mathcal B_\theta(Z)\).  With
\[
  \lambda_H:=e^{P_{Z_H}(\psi)},
\]
there exist \(h_H\in\mathcal B_\theta(Z)\), a non-zero positive
functional \(\eta_H\in\mathcal B_\theta(Z)^*\), and a bounded operator
\(\mathcal N_H\) on \(\mathcal B_\theta(Z)\) such that
\[
  \mathcal L_{\psi,H}h_H=\lambda_H h_H,\qquad
  \mathcal L_{\psi,H}^*\eta_H=\lambda_H\eta_H,\qquad
  \eta_H(h_H)=1,
\]
\[
  \mathcal L_{\psi,H}^n f
  =
  \lambda_H^n\eta_H(f)h_H+\mathcal N_H^n f,
  \qquad
  \rho(\mathcal N_H)<\lambda_H .
\]
Equivalently, for every \(f\in\mathcal B_\theta(Z)\) and every
\(\Lambda\in\mathcal B_\theta(Z)^*\),
\[
  \lambda_H^{-n}\Lambda(\mathcal L_{\psi,H}^n f)
  \longrightarrow
  \Lambda(h_H)\eta_H(f).
\]
Thus the leading open spectral data on the H\"older Banach space are
conclusions of the theorem, not hypotheses.
\end{theorem}

\begin{proof}
Write \(\abs{f}_\theta\) for the Lipschitz seminorm associated to
\(d_\theta\), and let \(p\) be a memory for the SFT \(Z\).  Boundedness is
the standard Ruelle estimate, and we recall the argument because the same
pairing is used below.  If \(x,x'\) agree for \(N\geq p\) symbols, then
every admissible safe one-step preimage of \(x\) is paired with the
preimage obtained by prepending the same safe symbol to \(x'\).  The
paired preimages agree for \(N+1\) symbols, and the Lipschitz property of
\(\psi\) gives
\[
  \abs{e^{\psi(y)}-e^{\psi(y')}}\leq C_\psi d_\theta(x,x').
\]
Since the number of one-step preimages is uniformly bounded,
\[
  \norm{\mathcal L_{\psi,H}f}_\theta
  \leq C\norm{f}_\theta
\]
for a constant \(C\) depending only on \(Z,\psi,H\), and \(\theta\).  The
finitely many cases \(N<p\) are absorbed by the sup-norm bound, since
\(d_\theta(x,x')\geq\theta^{p-1}\) there.

Let
\[
  \mathcal R:\mathcal B_\theta(Z)\longrightarrow
  \mathcal B_\theta(Z_H),
  \qquad
  \mathcal R f=f|_{Z_H},
\]
and let \(\mathcal E:=\ker\mathcal R\).  The restriction map is onto
because H\"older functions on the closed subset \(Z_H\) extend to
H\"older functions on the ambient symbolic metric space.  Explicitly, if
\(A\subset Z\) is closed and \(g:A\to\RR\) is \(L\)-Lipschitz, then
\[
  G(x):=\inf_{a\in A}\bigl(g(a)+Ld_\theta(x,a)\bigr)
\]
is an \(L\)-Lipschitz extension of \(g\) to \(Z\).  Applying this formula
to the real and imaginary parts of a complex-valued function proves the
surjectivity of \(\mathcal R\).  Moreover \(\mathcal E\) is a closed
invariant subspace and
\[
  \mathcal R\mathcal L_{\psi,H}
  =
  \mathcal L_{\psi|Z_H}\mathcal R,
\]
where \(\mathcal L_{\psi|Z_H}\) is the ordinary Ruelle operator on the
closed survivor subshift.  Indeed, if \(x\in Z_H\), every safe preimage
of \(x\) is again in \(Z_H\); a preimage from a transient safe symbol
would give that symbol an infinite safe continuation.  The classical
Ruelle--Perron--Frobenius theorem on the mixing SFT \(Z_H\) gives a
simple isolated eigenvalue
\(\lambda_H=e^{P_{Z_H}(\psi)}\) for
\(\mathcal L_{\psi|Z_H}\), with the rest of its spectrum contained in a
disk of radius \(r_H<\lambda_H\)
\citep{Bowen2008,Ruelle1978,Walters1975GMeasures}.

It remains to check that the part of the open operator acting on
\(\mathcal E\) has smaller spectral radius.  Let
\(\mathcal S_\infty\) be the survivor alphabet and
\(\mathcal T=\mathcal S\setminus\mathcal S_\infty\).  As in the proof of
Lemma~\ref{lem:holder-safe-word-pressure}, there is no safe edge from
\(\mathcal T\) to \(\mathcal S_\infty\), and the graph on
\(\mathcal T\) is acyclic.  Hence there is \(K\) such that every
sufficiently long safe word is a word over \(\mathcal S_\infty\) followed
by a transient suffix of length at most \(K\).

For later use set
\[
  W_n(x):=
  \sum_{\substack{\sigma^n y=x\\ y_0,\ldots,y_{n-1}\in\mathcal S}}
  e^{S_n\psi(y)} .
\]
The finite-suffix decomposition just recalled, bounded distortion, and
the Ruelle--Perron--Frobenius bound for
\(\mathcal L_{\psi|Z_H}\) give a constant \(C_W<\infty\) such that
\[
  W_n(x)\leq C_W\lambda_H^n
  \qquad (x\in Z,\ n\geq1).
\]
Indeed, after one of finitely many transient suffixes of length at most
\(K\) is fixed, the remaining sum is a survivor transfer-operator sum on
\(Z_H\), evaluated at a fixed survivor continuation and changed only by a
uniform distortion factor.

Take \(f\in\mathcal E\).  If \(y\) occurs in an \(n\)-step summand of
\(\mathcal L_{\psi,H}^n f(x)\), then \(y_0,\ldots,y_{n-1}\) are safe.
For \(n>K\), the first \(n-K\) symbols of \(y\) form a survivor word.
Choose \(\bar y\in Z_H\) with the same first \(n-K\) symbols.  Since
\(f|_{Z_H}=0\),
\[
  \abs{f(y)}
  =
  \abs{f(y)-f(\bar y)}
  \leq
  \norm{f}_\theta\,\theta^{n-K}.
\]
Together with the preceding bound on \(W_n(x)\), this gives
\[
  \norm{\mathcal L_{\psi,H}^n f}_\infty
  \leq
  C_2\theta^n e^{nP_{Z_H}(\psi)}\norm{f}_\theta .
\]

The H\"older seminorm satisfies the same bound.  Suppose first that
\(x,x'\) agree for \(N\geq p\) symbols, so \(d_\theta(x,x')=\theta^N\).
Pair the \(n\)-step safe preimages by the same safe prefix:
\(y=ux\) and \(y'=ux'\), \(\abs u=n\).  Then \(y\) and \(y'\) agree for
\(n+N\) symbols.  Hence
\[
  \abs{f(y)-f(y')}
  \leq
  \abs{f}_\theta\,\theta^{n+N}
  =
  \abs{f}_\theta\,\theta^n d_\theta(x,x').
\]
Also, by the previous comparison with \(Z_H\),
\[
  \abs{f(y')} \leq \norm{f}_\theta\theta^{n-K}.
\]
Finally the Bowen distortion estimate along the paired \(n\)-blocks gives
\[
  \abs{S_n\psi(y)-S_n\psi(y')}
  \leq
  C_\psi d_\theta(x,x'),
\]
and therefore
\[
  \abs{e^{S_n\psi(y)}-e^{S_n\psi(y')}}
  \leq
  C_\psi e^{S_n\psi(y)}d_\theta(x,x')
\]
after increasing \(C_\psi\).  For each paired summand,
\[
\begin{aligned}
 &\abs{e^{S_n\psi(y)}f(y)-e^{S_n\psi(y')}f(y')}        \\
 &\qquad\leq
  C e^{S_n\psi(y)}\norm{f}_\theta
  \theta^n d_\theta(x,x').
\end{aligned}
\]
Summing over safe prefixes and using \(W_n(x)\leq C_W\lambda_H^n\)
therefore yields
\[
  \abs{\mathcal L_{\psi,H}^n f}_\theta
  \leq
  C_3\theta^n e^{nP_{Z_H}(\psi)}\norm{f}_\theta .
\]
If \(N<p\), the same estimate follows from the sup-norm bound divided by
the lower bound \(d_\theta(x,x')\geq\theta^{p-1}\).  Thus the displayed
seminorm estimate holds for all \(x\neq x'\).
Thus
\[
  \rho\bigl(\mathcal L_{\psi,H}|_{\mathcal E}\bigr)
  \leq
  \theta e^{P_{Z_H}(\psi)}
  < \lambda_H .
\]

The invariant exact sequence
\[
  0\longrightarrow \mathcal E
  \longrightarrow \mathcal B_\theta(Z)
  \overset{\mathcal R}{\longrightarrow}
  \mathcal B_\theta(Z_H)
  \longrightarrow 0
\]
now gives the required spectral comparison.  We spell out the elementary
Banach-space argument.  Let \(T=\mathcal L_{\psi,H}\), let
\(T_{\mathcal E}=T|_{\mathcal E}\), and let \(\overline T\) be the
operator induced by \(T\) on the quotient
\(\mathcal B_\theta(Z)/\mathcal E\).  The map induced by
\(\mathcal R\) identifies \(\overline T\) with
\(\mathcal L_{\psi|Z_H}\).  If
\(\lambda\notin\spec(T_{\mathcal E})\cup\spec(\overline T)\), then
\(T-\lambda\) is invertible: solve first in the quotient, lift any
representative, and then correct the remaining error in \(\mathcal E\)
using \((T_{\mathcal E}-\lambda)^{-1}\).  This also gives a bounded
inverse, because the quotient inverse and
\((T_{\mathcal E}-\lambda)^{-1}\) are bounded.  Conversely, when
\(\lambda\notin\spec(T_{\mathcal E})\), invertibility of \(T-\lambda\)
would force invertibility of \(\overline T-\lambda\): the inverse of
\(T-\lambda\) preserves \(\mathcal E\), since the solution in
\(\mathcal E\) is unique.

Consequently, outside the disk of radius
\(\max\{\theta\lambda_H,r_H\}\), the spectrum of \(T\) is exactly the
corresponding spectrum of the survivor operator.  In particular
\(\lambda_H\) is an isolated spectral value of \(T\).  Since
\(\lambda_H\notin\spec(T_{\mathcal E})\) and the survivor eigenspace at
\(\lambda_H\) is one-dimensional, the same quotient-correction argument
lifts the survivor eigenfunction uniquely to a one-dimensional
eigenspace for \(T\).  The Riesz projection at \(\lambda_H\) is therefore
rank one.  If \(h_H\) spans its range and \(\eta_H^0\) is the positive
Ruelle--Perron--Frobenius eigenfunctional on
\(\mathcal B_\theta(Z_H)\), then
\[
  \eta_H(f):=\eta_H^0(\mathcal Rf)
\]
is a non-zero positive eigenfunctional for \(T^*\).  Normalizing
\(h_H\) so that \(\eta_H(h_H)=1\) and writing
\(\Pi_Hf=\eta_H(f)h_H\), the decomposition
\(\mathcal N_H:=T(I-\Pi_H)\) has
\(\rho(\mathcal N_H)<\lambda_H\), and the displayed spectral
decomposition follows.
\end{proof}

\begin{lemma}[One-step Gibbs error-survival comparison]
\label{lem:holder-one-step-error-comparison}
Let \(Z\) be a topologically mixing one-sided subshift of finite type,
let \(\psi\) be a H\"older potential, and let \(\nu_\psi\) be its
equilibrium state.  Let \(H\) be a non-empty proper one-step Markov hole
with safe alphabet \(\mathcal S\).  Assume that the survivor subshift
\[
  Z_H:=\{z\in Z:z_n\in\mathcal S\text{ for all }n\geq0\}
\]
is topologically mixing, and assume the finite two-terminal residue
reachability condition from
Lemma~\ref{lem:holder-extension-ambiguity}.  For
\[
  \tau_H(z):=\inf\{n\geq0:z_n\notin\mathcal S\},
  \qquad
  \mathsf P_R:=\{z:\tau_H(z)<\infty,\ \tau_H(z)\in R\bmod q\},
\]
there is \(\theta_0>0\) such that, for every \(n\geq1\),
\[
  \theta_0\,\nu_{\psi}(\tau_H\geq n)
  \leq
  \varepsilon_n(\mathsf P_R;\nu_{\psi})
  \leq
  \frac12\,\nu_{\psi}(\tau_H\geq n).
\]
Consequently
\[
  \lim_{n\to\infty}
  -\frac1n\log\varepsilon_n(\mathsf P_R;\nu_\psi)
  =
  P_Z(\psi)-P_{Z_H}(\psi).
\]
\end{lemma}

\begin{proof}
Lemma~\ref{lem:holder-extension-ambiguity} gives a constant \(c_0>0\),
uniform in the safe word \(w\), such that for every admissible safe
length-\(n\) word \(w\),
\[
  c_0
  \leq
  \nu_\psi(\mathsf P_R\mid[w])
  \leq
  1-c_0 .
\]
Hence
\[
  c_0
  \leq
  \min\{
    \nu_\psi(\mathsf P_R\mid[w]),
    1-\nu_\psi(\mathsf P_R\mid[w])
  \}
  \leq \frac12 .
\]
The safe-prefix decomposition
Lemma~\ref{lem:holder-safe-prefix-bayes-decomposition} therefore gives
\[
  \varepsilon_n(\mathsf P_R;\nu_\psi)
  =
  \sum_{\substack{w\in\mathcal L_n(Z)\\
                  w_0,\ldots,w_{n-1}\in\mathcal S}}
    \nu_\psi([w])\,
    \min\{
      \nu_\psi(\mathsf P_R\mid[w]),
      1-\nu_\psi(\mathsf P_R\mid[w])
    \}.
\]
The sum of the cylinder weights over safe words is
\(\nu_\psi(\tau_H\geq n)\).  Taking \(\theta_0=c_0\) proves the
two-sided comparison.

Lemma~\ref{lem:holder-safe-word-pressure} gives
\[
  \lim_{n\to\infty}
  -\frac1n\log\nu_\psi(\tau_H\geq n)
  =
  P_Z(\psi)-P_{Z_H}(\psi),
\]
and the comparison transfers the same exponential rate to
\(\varepsilon_n(\mathsf P_R;\nu_\psi)\).
\end{proof}

\begin{lemma}[Finite-cylinder block reduction]
\label{lem:holder-block-reduction}
Let \(Y\) be a topologically mixing one-sided subshift of finite type,
let \(\varphi\) be H\"older, and let \(\nu_\varphi\) be its equilibrium
state.  Suppose that \(H\subset Y\) is a finite union of cylinders and
choose \(\ell\geq1\) such that \(H\) is a union of \(\ell\)-cylinders.
Let
\[
  \Pi_\ell:Y\longrightarrow Y^{(\ell)}
\]
be the \(\ell\)-block conjugacy, let
\(\psi:=\varphi\circ\Pi_\ell^{-1}\), and let \(H^{(\ell)}\) be the
one-step hole in \(Y^{(\ell)}\) whose symbols are the \(\ell\)-blocks
contained in \(H\).  Then:
\begin{enumerate}[label=(\roman*)]
  \item \((\Pi_\ell)_\ast\nu_\varphi\) is the equilibrium state
    \(\nu_\psi\) of \(\psi\) on \(Y^{(\ell)}\);
  \item \(P_{Y^{(\ell)}}(\psi)=P_Y(\varphi)\), and the survivor subshift
    of \(H^{(\ell)}\) is conjugate to \(Y_H\), so its pressure is
    \(P_{Y_H}(\varphi)\);
  \item if \(m\geq\ell\) and \(n=m-\ell+1\), then the
    length-\(m\) cylinder \(\sigma\)-algebra on \(Y\) is the pullback of
    the length-\(n\) cylinder \(\sigma\)-algebra on \(Y^{(\ell)}\);
  \item first entry into \(H\) at time \(t\) is exactly first entry into
    \(H^{(\ell)}\) at block time \(t\).  Consequently first-entry residue
    events correspond under \(\Pi_\ell\), and their depth-\(m\) and
    depth-\(n\) Bayes errors are equal.
\end{enumerate}
\end{lemma}

\begin{proof}
The map \(\Pi_\ell\) sends \(y\) to the sequence of overlapping blocks
\((y_t\cdots y_{t+\ell-1})_{t\geq0}\).  This is a topological conjugacy
from \(Y\) to the usual \(\ell\)-block presentation.  The variational
principle is invariant under conjugacy, so the pushforward of the
equilibrium state of \(\varphi\) is the equilibrium state of
\(\psi=\varphi\circ\Pi_\ell^{-1}\), and the two ambient pressures agree.
The same conjugacy sends points that avoid \(H\) at every time precisely
to points whose block symbols avoid \(H^{(\ell)}\) at every time, giving
the pressure identity on the survivor subsystem.

A length-\(m\) word \(y_0\cdots y_{m-1}\), with \(m\geq\ell\), determines
exactly the \(n=m-\ell+1\) consecutive \(\ell\)-blocks
\[
  y_0\cdots y_{\ell-1},\;
  y_1\cdots y_\ell,\;\ldots,\;
  y_{m-\ell}\cdots y_{m-1},
\]
and the admissibility constraints in the block presentation recover the
original word from those blocks.  Thus the finite observation
\(\sigma\)-algebras coincide under pullback.  Finally,
\(\sigma^t y\in H\) if and only if the \(t\)-th block of
\(\Pi_\ell y\) belongs to \(H^{(\ell)}\).  The first-entry time and its
residue are therefore preserved, and the Bayes errors computed on the
corresponding finite \(\sigma\)-algebras are identical.
\end{proof}

\begin{theorem}[Block-coded H\"older operator and Gibbs escape comparison]
\label{thm:holder-finite-cylinder-escape}
Let \(Y\) be a topologically mixing one-sided subshift of finite type,
let \(\varphi\) be a H\"older potential, and let \(\nu_\varphi\) be its
equilibrium state.  Let \(H\subset Y\) be a non-empty proper finite union
of cylinders, and define
\[
  Y_H:=\{y\in Y:\sigma^n y\notin H\text{ for every }n\geq0\}.
\]
Fix \(q\geq 2\) and a non-empty proper subset
\(R\subsetneq \ZZ/q\ZZ\), and set
\[
  \tau_H(y):=\inf\{n\geq0:\sigma^n y\in H\},
  \qquad
  \mathsf P_R:=\{y:\tau_H(y)<\infty,\ \tau_H(y)\in R\bmod q\}.
\]
Choose \(\ell\geq 1\) so that \(H\) is a union of \(\ell\)-cylinders,
let \((Y^{(\ell)},\sigma^{(\ell)})\) be the \(\ell\)-block presentation
of \(Y\), let \(\psi\) be the lifted H\"older potential, and let
\(H^{(\ell)}\) be the corresponding one-step hole.  Let
\(\mathcal S\) be the safe block alphabet.

Assume the following two finite hypotheses.
\begin{enumerate}[label=(G\arabic*)]
  \item The survivor subshift of \(H^{(\ell)}\) is topologically mixing.
  \item In the directed graph with vertices
\[
  \mathcal S\times\ZZ/q\ZZ,\qquad \star_R,\quad\star_{R^c},
\]
with safe edges
\((i,r)\to(j,r+1)\) for admissible safe transitions \(i\to j\), and with
terminal edges \((i,r)\to\star_R\) or \((i,r)\to\star_{R^c}\) according
as an admissible next block enters the hole with residue \(r+1\) in
\(R\) or in \(R^c\), both terminal vertices are reachable from every
\((i,r)\in\mathcal S\times\ZZ/q\ZZ\).
\end{enumerate}

Choose a symbolic metric \(d_\theta\) on \(Y^{(\ell)}\) for which
\(\psi\) is Lipschitz, and let
\[
  (\mathcal L_{\psi,H^{(\ell)}}f)(x)
  :=
  \sum_{\substack{\sigma^{(\ell)} y=x\\ y_0\in\mathcal S}}
    e^{\psi(y)}f(y)
\]
be the open Ruelle operator on
\(\mathcal B_\theta(Y^{(\ell)})=C^\alpha(Y^{(\ell)})\).  Then
\(\mathcal L_{\psi,H^{(\ell)}}\) has the spectral decomposition from
Theorem~\ref{thm:holder-open-transfer}, with leading eigenvalue
\[
  \lambda_H
  =
  e^{P_{(Y^{(\ell)})_{H^{(\ell)}}}(\psi)}
  =
  e^{P_{Y_H}(\varphi)}.
\]

For the first-entry residue event \(\mathsf P_R^{(\ell)}\) in the block
system, let
\(\varepsilon_n^{(\ell)}(\mathsf P_R^{(\ell)};\nu_{\psi})\) denote the
Bayes error computed from block prefixes of length \(n\).  Then, for
every \(m\geq\ell\), with \(n=m-\ell+1\),
\[
  \varepsilon_m(\mathsf P_R;\nu_\varphi)
  =
  \varepsilon_n^{(\ell)}
    (\mathsf P_R^{(\ell)};\nu_{\psi}).
\]
There is \(\theta_0>0\), depending only on \(Y,\varphi,H,q,R\) and the
chosen block presentation, such that, for every \(n\geq1\),
\[
  \theta_0\,\nu_{\psi}(\tau_{H^{(\ell)}}\geq n)
  \leq
  \varepsilon_n^{(\ell)}
    (\mathsf P_R^{(\ell)};\nu_{\psi})
  \leq
  \frac12\,\nu_{\psi}(\tau_{H^{(\ell)}}\geq n).
\]
Consequently
\[
  \lim_{m\to\infty}
  -\frac1m\log
  \varepsilon_m(\mathsf P_R;\nu_\varphi)
  =
  P_Y(\varphi)-P_{Y_H}(\varphi).
\]
The operator assertion is the explicit \(C^\alpha\)-space statement
above applied to the block presentation.  The finite-observation
assertion is only the finite-prefix comparison and escape exponent; it
does not claim a Banach-space cyclotomic resolvent.
\end{theorem}

\begin{proof}
The equality of Bayes errors and the pressure identifications under
block recoding are exactly Lemma~\ref{lem:holder-block-reduction}.  The
operator statement is Theorem~\ref{thm:holder-open-transfer} applied to
the block system \(Y^{(\ell)}\), the H\"older potential \(\psi\), and the
one-step hole \(H^{(\ell)}\); the block pressure identity in
Lemma~\ref{lem:holder-block-reduction} gives the displayed value of
\(\lambda_H\).

Apply Lemma~\ref{lem:holder-one-step-error-comparison} to the block
system \(Y^{(\ell)}\), the potential \(\psi\), and the one-step hole
\(H^{(\ell)}\).  Its two-sided comparison is exactly the displayed
comparison for
\(\varepsilon_n^{(\ell)}(\mathsf P_R^{(\ell)};\nu_\psi)\), and it also
gives
\[
  \lim_{n\to\infty}
  -\frac1n\log
  \varepsilon_n^{(\ell)}(\mathsf P_R^{(\ell)};\nu_\psi)
  =
  P_{Y^{(\ell)}}(\psi)
  -P_{(Y^{(\ell)})_{H^{(\ell)}}}(\psi).
\]
The fixed shift \(n=m-\ell+1\) does not affect the exponential rate.
Block conjugacy preserves pressure and identifies
\((Y^{(\ell)})_{H^{(\ell)}}\) with \(Y_H\).  Hence the last displayed
quantity is \(P_Y(\varphi)-P_{Y_H}(\varphi)\), and the block-recoding
equality gives the claimed limit for
\(\varepsilon_m(\mathsf P_R;\nu_\varphi)\).
\end{proof}
